\label{section:GSTR}
The \mfus\ \textsf{GSTR} package assigns recharge that varies in space and time using ArcASCII raster snapshots.
It operates at the timestep level, independently of the existing \textsf{RCH} and \textsf{RTS} packages.
Existing \textsf{RCH}/\textsf{RTS} input files remain fully supported; when both \textsf{GSTR} and \textsf{RCH} (or \textsf{RTS}) are active, recharge fluxes from each package are added to the flow equation.

\subsection{MUT instructions}
Select the cells that belong to the instance, optionally name the instance, then assign snapshots on the target domain.

\ins{gstr instance name}
{
    \squish
    \begin{enumerate}
    \item \str{name} Budget label for the instance (at most 16 characters).
    \end{enumerate}
    The name is applied to the next \texttt{gwf gstr}, \texttt{swf gstr}, or \texttt{cln gstr} instruction.
    If omitted, MUT assigns a default name \texttt{GSTR\_$n$}.
    \squish
}

\ins{gwf gstr}
{
    \squish
    \begin{enumerate}
    \item \rnum{FAC} Multiplier on sampled rates (use $1$ for none).
    \item One or more lines \rnum{time} \str{raster.asc} (ascending times, model time units).
    \item \texttt{end gstr}
    \end{enumerate}
    Currently chosen \gwf\ cells are written with global node number and centroid $(x,y)$.
    Identical forms exist for \swf\ (\texttt{swf gstr}, IDOMAIN$=3$) and \cln\ (\texttt{cln gstr}, IDOMAIN$=2$).
    \squish
}

Example (SWF rainfall with start and end rasters):
\begin{verbatim}
active domain
swf
    choose all cells
    gstr instance name
    Rain_SWF
    swf gstr
    1.0
    0.0    rain_t0.asc
    20.0   rain_t20.asc
    end gstr
\end{verbatim}

MUT writes \texttt{\textit{prefix}.gstr} and a name-file entry \texttt{GSTR}.
A KAERI/KURT demonstration lives under
\texttt{KURT\_Model/2\_models/2\_TestTransientRainfall/4\_GSTR}.

\subsection{Name file}
Add one name-file entry:
\begin{verbatim}
GSTR  <unit>  <prefix>.gstr
\end{verbatim}

\subsection{Input file format}
Comment lines may begin with \texttt{\#} or \texttt{!}. Blank lines are skipped.

\begin{verbatim}
NGSTR  IGSTRCB
INSTANCE  <name>
  IDOMAIN  NCELLS  NSNAPS  [FAC]
  <inode>  <x>  <y>
  ...
  <time>   <raster.asc>
  ...
INSTANCE  <name2>
  ...
\end{verbatim}

\begin{description}
\item[\texttt{NGSTR}] Number of named GSTR instances.
\item[\texttt{IGSTRCB}] Unit number for cell-by-cell budget terms (\texttt{0} = none).
\item[\texttt{INSTANCE} \textit{name}] Starts a named instance. The name (up to 16 characters) is used as the volumetric budget label for that instance.
\item[\texttt{IDOMAIN}] Domain code: \texttt{1}=\gwf, \texttt{2}=\cln, \texttt{3}=\swf. Used for validation; recharge is applied at the listed global node numbers.
\item[\texttt{NCELLS}] Number of cells in the instance (all domain cells or a subset).
\item[\texttt{NSNAPS}] Number of time snapshots (must be $\ge 1$).
\item[\texttt{FAC}] Optional multiplier applied to sampled recharge rates (default $1$).
\item[\texttt{inode x y}] Global node number and cell coordinates used to sample rasters. Coordinates must be in the same length units and coordinate system as the ArcASCII grids.
\item[\texttt{time raster.asc}] Snapshot time (model time units, ascending) and path to an ArcASCII recharge raster ($L/T$).
\end{description}

\subsection{ArcASCII rasters}
Each snapshot file is a standard ArcASCII grid with header keywords
\texttt{ncols}, \texttt{nrows}, \texttt{xllcorner} or \texttt{xllcenter},
\texttt{yllcorner} or \texttt{yllcenter}, \texttt{cellsize}, and optional \texttt{nodata\_value},
followed by row-major cell values (north to south).
Bilinear interpolation is used at each cell $(x,y)$.
Samples outside the raster extent, or involving a nodata corner, are treated as zero recharge.

\subsection{Time interpolation}
At each timestep, \mfus\ uses elapsed simulation time (\texttt{TOTIM}) to select the raster snapshots immediately before and after the current time and linearly interpolates the sampled rate for each cell.
Outside the snapshot range, the nearest endpoint snapshot is held (no extrapolation beyond the first or last raster).
The volumetric flux applied to a node is
\[
Q = \mathrm{FAC}\times r(x,y,t)\times A_n
\]
where $r$ is the interpolated rate ($L/T$) and $A_n$ is the node area stored by \mfus\ (GWF area, SWF face area, or CLN conduit area as appropriate).

\subsection{Budget}
Each named instance contributes its own budget term using the instance name as the label.
\textsf{GSTR} does not replace or modify \textsf{RCH}/\textsf{RTS} budget terms.

\subsection{Post-processing to Tecplot}
After \texttt{postprocess existing modflow model}, MUT reads the \textsf{GSTR} package and reconstructs applied rates by sampling the snapshot rasters (and a mid-interval interpolated zone when two or more snaps exist). Outputs include:
\begin{verbatim}
_posto.o.<prefix>.SWF_GSTR.tecplot.dat
_posto.o.<prefix>.SWF_GSTR_TS.tecplot.dat
\end{verbatim}
(and analogous \texttt{GWF\_GSTR} / \texttt{CLN\_GSTR} files by domain).
Named instance totals also appear in the listing-file volume-budget Tecplot export.
